\documentclass[10pt,a4paper]{amsart}
\usepackage[margin=19mm]{geometry}
\usepackage{amsmath,amssymb,mathtools}
\usepackage[scr=rsfs,cal=euler]{mathalfa}
\usepackage{xfrac}
\usepackage[unicode,hidelinks,bookmarks=false]{hyperref}
\newcommand{\ie}{i.e.\ }
\newcommand{\cf}{cf.\ }
\newcommand{\resp}{respectively\ }
\newcommand{\Subsection}{\subsection}
\providecommand{\nobreakdash}{\nobreak\hbox{-}}
\setlength{\emergencystretch}{2em}
\multlinegap0pt
\usepackage[shortlabels]{enumitem}
\usepackage{amssymb}
\usepackage{bbm}
\usepackage{mathtools}
\usepackage{tikz}
\newtheorem{definition}{Definition}
\newtheorem{lemma}{Lemma}
\newcommand{\R}{\mathbb{R}}
\newcommand{\acts}{\curvearrowright}
\renewcommand{\epsilon}{\varepsilon}
\DeclareMathOperator{\inter}{\mathrm{int}}
\newcommand{\mfR}{\mathfrak{R}}
\DeclareMathOperator{\proj}{\mathrm{proj}}
\DeclareMathOperator{\ran}{\mathrm{ran}}
\newcommand{\unifm}{\mathsf{D}}
\newcommand{\vietoris}[2][]{\mathcal{K}_{#1}(#2)}
\title{\vspace{-2.2cm} Equivariant Borel liftings in complex analysis and PDE}
\author{Konstantin Slutsky, Mikhail Sodin, Aron Wennman}
\date{Source: arXiv:2507.12058v2 (CC BY 4.0)}
\begin{document}
\maketitle

\noindent\emph{Source and licence.} \vspace{-2.2cm} Equivariant Borel liftings in complex analysis and PDE. Version arXiv:2507.12058v2 (CC BY 4.0), distributed by arXiv under the Creative Commons Attribution 4.0 International licence, http://creativecommons.org/licenses/by/4.0/, as recorded in the arXiv licence metadata for this exact version and shown on the abstract page https://arxiv.org/abs/2507.12058v2. Full licence notice: \texttt{licenses/arxiv-2507.12058-CC-BY-4.0.txt}.\par\smallskip

\noindent\textit{Editorial scope.} Counted unit: the article's lemma lem:runge-equivariant-uniformization (source0.tex lines 1284--1289). The article's complete original proof of it is delivered with it (source0.tex lines 1292--1424, one \texttt{proof} environment of 133 lines). No mathematics is written, repaired or re-derived: the statement and the proof are the article's own lines, in the article's own notation, and no retained line is substituted or re-flowed.  Retained uncounted as the source-local context the proof needs: lines 972--1003; lines 1058--1063; lines 1210--1230.  Nothing was omitted from any retained span, so the package carries no omission record.  No retained line contains a citation, so the wrapper carries no bibliography and no bibliography entry.  No retained line contains a figure environment, an image-inclusion command or drawing code, so no image is delivered and the package declares no asset dependency.  Editorial formatting only, in addition: an AMS \texttt{amsart} preamble that restates the article's own declarations and macros used by the retained lines, namely the environments \texttt{definition}, \texttt{lemma} on separate counters, together with the standard packages those lines need.  The retained environments print in this document as Definition 1, Lemma 1 in order of appearance, whereas the article prints the same items with the numbering of its own full document.  The complete native source of the article as distributed in the arXiv e-print for this exact version is bundled unchanged as \texttt{sources/181667/source0.tex}.
\begin{definition}
  \label{def:toast}
  Let \(a_{X} : G \acts X\) be a free Borel action of a locally compact Polish group on a standard
  Borel space. A Borel \emph{toast} for \(a_{X}\) is a sequence
  \((\mathcal{C}_{n}, \lambda_{n})_{n}\) of Borel sets \(\mathcal{C}_{n} \subseteq X\) and Borel
  functions \(\lambda : \mathcal{C}_{n} \to \vietoris[*]{G}\) satisfying the following
  conditions. For each \(n\), define \(R_{n}(c) = \lambda_{n}(c)\cdot c\) for
  \(c \in \mathcal{C}_{n}\), and let \(X_{n} = \bigcup_{c \in \mathcal{C}_{n}} R_{n}(c)\). Then:
  \begin{enumerate}[leftmargin=1cm, label={\sf \arabic*}., ref={\sf \arabic*}]
  \item\label{item:toast-disjoint} \(R_{n}(c_{n}) \cap R_{n}(c'_{n}) = \varnothing\) for all
    distinct \(c_{n}, c'_{n} \in \mathcal{C}_{n}\).
  \item\label{item:toast-coherent} For all \(m < n\), \(c_{m} \in \mathcal{C}_{m}\), and
    \(c_{n} \in \mathcal{C}_{n}\), either \(R_{m}(c_{m}) \cap R_{n}(c_{n}) = \varnothing\) or
    \(R_{m}(c_{m}) \subseteq R_{n}(c_{n})\).
  \item\label{item:toast-layered} For all \(c_{n} \in \mathcal{C}_{n}\) there exists
    \(c_{n+1} \in \mathcal{C}_{n+1}\) such that \(R_{n}(c_{n}) \subseteq R_{n+1}(c_{n+1})\).
  \item\label{item:toast-directed} For all \(c_{m_{1}} \in \mathcal{C}_{m_{1}}\) and
    \(c_{m_{2}} \in \mathcal{C}_{m_{2}}\), there exist \(n > m_{1}, m_{2}\) and an element
    \(c_{n} \in \mathcal{C}_{n}\) such that \(R_{m_{i}}(c_{m_{i}}) \subseteq \inter R_{n}(c_{n})\)
    for \(i = 1,2\), where \(\inter R_{n}(c_{n}) = (\inter \lambda_{n}(c_{n})) \cdot c_{n}\).
  \item\label{item:lacunary} There exists a neighborhood of the identity \(U \subseteq G\) such that
    \(U \subseteq \lambda_{n}(c_{n})\) for all \(c_{n} \in \mathcal{C}_{n}\) and all \(n\).
  \item\label{item:toast-exhaustive} \(\bigcup_{n} X_{n} = X\).
  \item\label{item:toast-rational} The range \(\ran \lambda_{n}\) is countable for each \(n\) and
    \[\{\rho(c_{m},c_{n}) : c_{m} \in \mathcal{C}_{m}, c_{n} \in \mathcal{C}_{n}, c_{m} E_{a_{X}}
      c_{n}\}\]
    is countable, where \(\rho : E_{a_{X}} \to G\) is the cocycle defined by the condition
    \(\rho(x_{1},x_{2})x_{1} = x_{2}\).
  \end{enumerate}
  We say that \((\mathcal{C}_{n}, \lambda_{n})_{n}\) is an \(\mfR\)-toast, where
  \(\mfR \subseteq \vietoris{G}\), if \(\ran \lambda_{n} \subseteq \mfR\) for all \(n\).
\end{definition}

\begin{lemma}
  \label{lem:exhaustive-taust}
  Let \((\mathcal{C}_{n},\lambda_{n})_{n}\) be a Borel toast for \(G \acts X\).  For all \(x \in X\)
  and \(K \in \vietoris{G}\), there exist \(n\) and \(c_{n} \in \mathcal{C}_{n}\) such that
  \(K \cdot x \subseteq \inter R_{n}(c_{n})\).
\end{lemma}

\begin{definition}
  \label{def:U-family}
  Let \(a_{Z} : G \acts Z\) be a Borel \(G\)-action.  A family
  \(\unifm = (d_{K})_{K \in \vietoris{G}}\) of pseudometrics on \(Z\) is said to be an
  \emph{\(a_{Z}\)-family} if it satisfies the following conditions:
  \begin{enumerate}[leftmargin=1cm, label={\sf \arabic*}., ref={\sf \arabic*}]
  \item\label{item:unif-monotone} \(d_{K}(z_{1},z_{2}) \le d_{K'}(z_{1}, z_{2})\) for all
    \(z_{1}, z_{2} \in Z\) and \(K, K' \in \vietoris{G}\) satisfying \(K \subseteq K'\).
  \item\label{item:unif-shift} \(d_{K}(gz_{1}, gz_{2}) = d_{Kg}(z_{1},z_{2})\) for all
    \(z_{1},z_{2} \in Z\), \(g \in G\), and \(K \in \vietoris{G}\).
  \item\label{item:unif-hausdorff} The family \(\unifm\) separates points: for any distinct
    \(z_{1}, z_{2} \in Z\) there exists \(K \in \vietoris{G}\) such that
    \(d_{K}(z_{1}, z_{2}) > 0\).
  \item\label{item:unif-complete} The uniformity generated by \( \unifm \) is complete: if
    \((z_{n})_{n}\) is \(\unifm\)-Cauchy (i.e., \(d_{K}\)-Cauchy for every \(K \in \vietoris{G}\)),
    then there exists some \(z \in Z\) such that \(d_{K}(z_{n},z) \to 0\) for all
    \(K \in \vietoris{G}\).
  \end{enumerate}
  An \(a_{Z}\)-family is \emph{Borel} if each \(d_{K}\) is Borel as a map
  \(d_{K} : Z \times Z \to \R^{\ge 0} \cup \{\infty\}\).
\end{definition}

\begin{lemma}
  \label{lem:runge-equivariant-uniformization}
  Suppose that \(\unifm\) satisfies the \((\mfR, P)\)-Runge property and assume that \(a_{X}\)
  admits a Borel \(\mfR\)-toast.  If each fiber \(P_{x}\), \(x \in \proj_{X}(P)\), is
  \(\unifm\)-closed, then \(P\) admits an \((a_{X}, a_{Z})\)-equivariant Borel uniformization.
\end{lemma}

\begin{proof}
  By the assumption, the set \(P\) admits a Borel uniformization, which implies that
  \(\proj_{X}(P)\) is a Borel subset of \(X\). Additionally, \(\proj_{X}(P)\) is
  \(a_{X}\)-invariant. Without loss of generality, we may substitute \(\proj_{X}(P)\) for \(X\),
  thereby assuming that \(\proj_{X}(P) = X\).

  Recall that \(a_{X}\) is assumed to be free, which allows us to define the cocycle map
  \(\rho : E_{G}^{X} \to G\) by the condition \(\rho(x, y)x = y\) for \(x E_{G}^{X} y\).  Let
  \((\mathcal{C}_{n}, \lambda_{n})_{n}\) be a Borel \(\mfR\)-toast for \(a_{X}\).

  We inductively construct a sequence of Borel maps \(\xi_{n} : \mathcal{C}_{n} \to Z\) such that
  \((c_{n}, \xi_{n}(c_{n})) \in P\) and the following condition holds:
  \begin{equation}
    \label{eq:xi-condition}
    d_{\lambda_{n-1}(c_{n-1})\rho(c_{n},c_{n-1})}(\xi_{n}(c_{n}),
    \rho(c_{n-1},c_{n})\cdot \xi_{n-1}(c_{n-1})) < 2^{-n}
  \end{equation}
  for all \(c_{n-1} \in \mathcal{C}_{n-1}\) and \(c_{n} \in \mathcal{C}_{n}\) such that \(c_{n-1} \)
  is a predecessor of \(c_{n}\), i.e., \(R_{n-1}(c_{n-1}) \subseteq R_{n}(c_{n})\).

  For the base case, define \(\xi_{0} : \mathcal{C}_{0} \to Z\) to be the restriction of any Borel
  uniformization of \(P\) to \(\mathcal{C}_{0}\). Now, assume that \(\xi_{n-1}\) has already been
  constructed. Choose an element \(c_{n} \in \mathcal{C}_{n}\), and let
  \(c_{n-1}^{1}, \ldots, c_{n-1}^{m} \in \mathcal{C}_{n-1}\) denote all the elements of
  \(\mathcal{C}_{n-1}\) satisfying \(R_{n-1}(c_{n-1}^{i}) \subseteq R_{n}(c_{n})\).  Define compact
  sets \(K_{i} = \lambda_{n-1}(c_{n-1}^{i})\rho(c_{n},c_{n-1}^{i})\) for \(1 \le i \le m\). By the
  invariance of \(\mfR\), each \(K_{i}\) lies in \(\mfR\), and the definition of the toast ensures
  that the sets \(K_{i}\), \(1 \le i \le m\), are pairwise disjoint.

  Let \(f : P^{(m)} \to Z\) be the map provided by the \((\mfR, P)\)-Runge property for the sets
  \(K_{1}, \ldots, K_{m}\) and \(\epsilon = 2^{-n}\). Observe that, for each \(1 \le i \le m\),
  \((c_{n-1}^{i}, \xi_{n-1}(c_{n-1}^{i})) \in P\) implies
  \((c_{n}, \rho(c_{n-1}^{i},c_{n})\cdot \xi_{n-1}(c_{n-1}^{i})) \in P\). We can thus define
  \[
    \xi_{n}(c_{n}) = f(c_{n}, \rho(c_{n-1}^{1},c_{n})\cdot \xi_{n-1}(c_{n-1}^{1}), \ldots,
    \rho(c_{n-1}^{m},c_{n})\cdot \xi_{n-1}(c_{n-1}^{m})).
  \]
  The defining property of \(f\) ensures that
  \[
    d_{K_{i}}(\xi_{n}(c_{n}), \rho(c_{n-1}^{i},c_{n})\cdot \xi_{n-1}(c_{n-1}^{i})) < 2^{-n}
  \]
  for all \(1 \le i \le m\).  Thus, \(\xi_{n}\) satisfies Eq.~\eqref{eq:xi-condition}.  Note that
  item~\eqref{item:toast-rational} of the definition of a Borel toast ensures that there are only
  countably many distinct possibilities for the sets \(K_{i}\), which allows us to perform the
  construction above in a Borel way over all \(\mathcal{C}_{n}\).

  Let \(X_{n} = \bigcup_{c_{n} \in \mathcal{C}_{n}} R_{n}(c_{n})\), and define functions
  \(\beta_{n} : X_{n} \to \mathcal{C}_{n}\) by the condition \(x \in R_{n}(\beta_{n}(x))\) for all
  \(x \in X_{n}\). Next, define functions \(\psi_{n} : X_{n} \to Z\) as
  \begin{equation}
    \label{eq:psi-n-definition}
      \psi_{n}(x) = \rho(\beta_{n}(x), x)\cdot \xi_{n}(\beta_{n}(x)).
  \end{equation}
  Since \((c_{n}, \xi_{n}(c_{n})) \in P\) for all \(c_{n} \in \mathcal{C}_{n}\) and using the
  \((a_{X},a_{Z})\)-equivariance of \(P\), it follows that
  \[
    P \ni \big(\rho(\beta_{n}(x), x)\cdot\beta_{n}(x), \rho(\beta_{n}(x),
    x)\cdot\xi_{n}(\beta_{n}(x))\big) = (x, \psi_{n}(x)).
  \]
  Thus, each \(\psi_{n}\) serves as a Borel uniformization of \(P \cap (X_{n} \times Z)\).

  We claim that the sequence \((\psi_{n}(x))_{n}\) is \(\unifm\)-Cauchy for every \(x \in X\). To
  verify this, it is enough to show that for each \(K \in \vietoris{G}\), every \(x \in X\), and all
  sufficiently large \(n\), the inequality \(d_{K}(\psi_{n}(x), \psi_{n-1}(x)) < 2^{-n}\) holds. Fix
  \(K \in \vietoris{G}\) and define \(c_{n-1} = \beta_{n-1}(x)\) and \(c_{n} = \beta_{n}(x)\). For
  sufficiently large \(n\), Lemma~\ref{lem:exhaustive-taust} gives
  \[
    K\cdot x \subseteq \lambda_{n-1}(c_{n-1})\cdot c_{n-1} \subseteq \lambda_{n}(c_{n})\cdot c_{n}.
  \]
  Since the action \(a_{X}\) is free, it follows that
  \begin{equation}
    \label{eq:containment}
    K\rho(c_{n},x) = K\rho(c_{n-1},x)\rho(c_{n},c_{n-1}) \subseteq
    \lambda_{n-1}(c_{n-1})\rho(c_{n},c_{n-1}).
  \end{equation}
  We now have the following equalities
  \begin{displaymath}
    \begin{aligned}
      d_{K}(\psi_{n}(x), \psi_{n-1}(x))
      &= d_{K}(\rho(c_{n},x)\cdot\xi_{n}(c_{n}), \rho(c_{n-1},x)\cdot\xi_{n-1}(c_{n-1})) \\
      &= d_{K\rho(c_{n}, x)}(\xi_{n}(c_{n}), \rho(x,c_{n})\rho(c_{n-1},x)\cdot\xi_{n-1}(c_{n-1})) \\
      &= d_{K\rho(c_{n}, x)}(\xi_{n}(c_{n}), \rho(c_{n-1},c_{n})\cdot\xi_{n-1}(c_{n-1})) \\
    \end{aligned}
  \end{displaymath}
  which yields the estimate
  \begin{equation}
    \label{eq:psi-n-cauchy}
    \begin{aligned}
      d_{K}(\psi_{n}(x), \psi_{n-1}(x))
      & = d_{K\rho(c_{n}, x)}(\xi_{n}(c_{n}),
        \rho(c_{n-1},c_{n})\cdot\xi_{n-1}(c_{n-1})) \\
      &\le
        d_{\lambda_{n-1}(c_{n-1})\rho(c_{n},c_{n-1})}(\xi_{n}(c_{n}),
        \rho(c_{n-1},c_{n})\cdot\xi_{n-1}(c_{n-1})) \\
      &< 2^{-n},
    \end{aligned}
  \end{equation}
  where the first inequality is due to Eq.~\eqref{eq:containment} and the second follows from
  Eq.~\eqref{eq:xi-condition}.

  We have established that \((\psi_{n}(x))_{n}\) is \(\unifm\)-Cauchy for every \(x \in X\). By
  item~\eqref{item:unif-complete} of Definition~\ref{def:U-family}, the family of pseudometrics
  \(\unifm\) is complete. Consequently, for each \(x \in X\), the limit
  \(\psi(x) = \textrm{\(\unifm\)-}\lim_{n}\psi_{n}(x)\) exists and is unique, as \(\unifm\)
  separates points. Furthermore, \(\psi\) is a Borel map because \(\unifm\) is Borel. To verify
  this, take an exhaustion \((K_{n})_{n}\) of \(G\) by compact sets. The graph of \(\psi\) can be
  written as
  \[
    \Bigl\{(x,y) : \forall k\ \forall m\ \exists N\ (\forall n\ge N)\ [d_{K_{m}}(\psi_{n}(x),y) <
    1/k]\Bigr\},
  \]
  which is Borel since the pseudometrics \(d_{K_{m}}\) and the functions \(\psi_{n}\) are Borel.

  We claim that \(\psi\) is \((a_{X}, a_{Z})\)-equivariant.  Observe that
  \(\rho(y,g\cdot x) = g \rho(y,x)\) for all \(x,y \in X\) and \(g \in G\).  Moreover, for all
  sufficiently large \(n\), we have \(\beta_{n}(g\cdot x) = \beta_{n}(x)\) and consequently
  \begin{displaymath}
    \psi_{n}(g\cdot x) = \rho(\beta_{n}(g\cdot x), g\cdot x) \cdot \xi_{n}(\beta_{n}(g\cdot x)) =
    g\rho(\beta_{n}(x), x)\cdot \xi_{n}(\beta_{n}(x)) = g\cdot \psi_{n}(x).
  \end{displaymath}
  Taking the limit as \(n \to \infty\), we obtain
  \[\psi(g\cdot x) = \unifm\textrm{-}\lim_{n}\psi_{n}(g\cdot x)
    = \unifm\textrm{-}\lim_{n}g\cdot \psi_{n}(x) = g\cdot (\unifm\textrm{-}\lim_{n}\psi_{n}(x)) =
    g\cdot \psi(x),\]
  where the penultimate equality follows from the continuity of the map \(z \mapsto g\cdot z\) in
  the topology of \(\unifm\), given by item~\eqref{item:unif-shift} of
  Definition~\ref{def:U-family}.  Thus, \(\psi \) is \(G\)-equivariant.

  Recall that each \(\psi_{n}\) is a Borel uniformization of \(P \cap (X_{n} \times Z)\).  Since
  slices \(P_{x}\) are assumed to be \(\unifm\)-closed, we conclude that \((x,\psi(x)) \in P\) holds
  for all \( x \in X\) and \(\psi\) is therefore a Borel \((a_{X}, a_{Z})\)-equivariant
  uniformization of \(P\).
\end{proof}

\end{document}
